\documentclass[11pt,a4paper,UTF8]{ctexart}
\usepackage{geometry}
\geometry{left=18mm,right=18mm,top=24mm,bottom=24mm,headheight=22pt,footskip=12mm}
\usepackage{amsmath,amssymb,mathtools,bm,esint,extarrows}
\usepackage{xcolor}
\usepackage{tcolorbox}
\tcbuselibrary{skins,breakable}
\usepackage{fancyhdr}
\usepackage{titlesec}
\usepackage{hyperref}
\usepackage{tikz}
\usepackage{booktabs}
\usepackage{tabularx}

% --- Color Palette (Purple & DeepPink Academic Theme) ---
\definecolor{DeepPurple}{HTML}{4C1D95}
\definecolor{TitlePurple}{HTML}{581C87}
\definecolor{SubPurple}{HTML}{6B21A8}
\definecolor{DeepPink}{HTML}{FF1493}
\definecolor{LightPinkBg}{HTML}{FFF5F8}
\definecolor{LightPurpleBg}{HTML}{F8F5FF}
\definecolor{GoldAccent}{HTML}{D97706}
\definecolor{DarkText}{HTML}{1F2937}

\hypersetup{
    colorlinks=true,
    linkcolor=TitlePurple,
    citecolor=DeepPink,
    urlcolor=SubPurple,
    pdfborder={0 0 0}
}

% --- Typography & Spacing ---
\linespread{1.3}
\setlength{\parskip}{0.35em plus 0.1em minus 0.1em}
\setlength{\parindent}{0pt}

% --- Section Titles Styling ---
\titleformat{\section}
  {\Large\bfseries\color{TitlePurple}}{\thesection}{1em}{}
  [\vspace{1mm}{\color{TitlePurple!30}\hrule height 1pt}]
\titleformat{\subsection}
  {\large\bfseries\color{SubPurple}}{\thesubsection}{0.8em}{}
\titleformat{\subsubsection}
  {\normalsize\bfseries\color{DeepPurple}}{\thesubsubsection}{0.6em}{}

% --- tcolorbox Environments ---
% 1. Theorem / Formula / Method / Analysis Box (Deep Purple)
\newtcolorbox{thmbox}[1][]{
    enhanced,
    breakable,
    colback=LightPurpleBg,
    colframe=TitlePurple,
    arc=3pt,
    boxrule=0pt,
    leftrule=3.5pt,
    left=10pt,
    right=10pt,
    top=8pt,
    bottom=8pt,
    title=#1,
    coltitle=white,
    colbacktitle=TitlePurple,
    fonttitle=\bfseries\small,
    attach boxed title to top left={yshift=-2mm, xshift=4mm},
    boxed title style={boxrule=0pt, arc=2pt}
}

% 2. Solution / Formula / Derivation Box (DeepPink)
\newtcolorbox{solbox}[1][]{
    enhanced,
    breakable,
    colback=LightPinkBg,
    colframe=DeepPink,
    arc=3pt,
    boxrule=0pt,
    leftrule=3.5pt,
    left=10pt,
    right=10pt,
    top=8pt,
    bottom=8pt,
    title=#1,
    coltitle=white,
    colbacktitle=DeepPink,
    fonttitle=\bfseries\small,
    attach boxed title to top left={yshift=-2mm, xshift=4mm},
    boxed title style={boxrule=0pt, arc=2pt}
}

% 3. Learning Goals / Summary Box
\newtcolorbox{targetbox}[1][]{
    enhanced,
    breakable,
    colback=LightPurpleBg,
    colframe=DeepPurple,
    arc=3pt,
    boxrule=1pt,
    left=10pt,
    right=10pt,
    top=8pt,
    bottom=8pt,
    title=#1,
    coltitle=white,
    colbacktitle=DeepPurple,
    fonttitle=\bfseries\small,
    attach boxed title to top left={yshift=-2mm, xshift=4mm},
    boxed title style={boxrule=0pt, arc=2pt}
}

\begin{document}

% ------------------- 讲义封面 (Cover Page) -------------------
\begin{titlepage}
\thispagestyle{empty}
\begin{tikzpicture}[remember picture,overlay]
    \fill[DeepPurple] (current page.north west) rectangle ([yshift=-7cm]current page.north east);
    \draw[DeepPink, line width=3.5pt] ([yshift=-7cm]current page.north west) -- ([yshift=-7cm]current page.north east);
    \draw[GoldAccent, line width=1pt] ([yshift=-7.12cm]current page.north west) -- ([yshift=-7.12cm]current page.north east);
    
    \node[anchor=north east, opacity=0.12, text=white] at ([xshift=-1.5cm, yshift=-0.5cm]current page.north east) {
        \fontsize{70}{70}\selectfont\bfseries 2027
    };
\end{tikzpicture}

\vspace*{0.8cm}
\begin{center}
    {\color{white}\fontsize{26}{32}\selectfont\bfseries 2027 考研数学(二) 核心讲义}\\[0.6cm]
    {\color{white}\fontsize{13}{18}\selectfont\bfseries 全国硕士研究生招生考试 · 统考数学(二) 核心精讲全书}\\[1.5cm]
\end{center}

\vspace{3.5cm}
\begin{center}
    {\color{GoldAccent}\large\bfseries $\blacktriangleright$ 基础强化 · 核心题解 · 考点深水区 $\blacktriangleleft$}\\[0.8cm]
    {\color{TitlePurple}\fontsize{24}{28}\selectfont\bfseries [强化]一元积分学的应用(物理和经济)}\\[0.6cm]
    {\color{DeepPink}\large\bfseries —— 考研高阶冲刺与深水区专攻 ——}\\[1.5cm]
\end{center}

\vfill

\begin{center}
\begin{tcolorbox}[
    colback=white,
    colframe=DeepPurple,
    width=0.88\textwidth,
    arc=4pt,
    boxrule=1.2pt,
    halign=center,
    drop shadow=gray!30
]
    \vspace{3mm}
    {\bfseries\large 编著团队：EscoffierZhou}\\[2.5mm]
    {\bfseries\color{TitlePurple} 目标院校：中国科学院大学 (UCAS) 计算机专硕 (22408)}\\[2.5mm]
    {\small\color{gray} 备考铁律：手算真题数字 · 错题白纸重做 · 408 客观题为王}\\[2.5mm]
    {\small 编制版本：2027 届首发版 · \today}
    \vspace{2mm}
\end{tcolorbox}
\end{center}
\vspace{0.8cm}
\end{titlepage}

% ------------------- 目录与页面主体配置 -------------------
\pagestyle{fancy}
\fancyhf{}
\fancyhead[L]{\small\color{gray}2027 考研数学(二) · 核心讲义系列}
\fancyhead[R]{\small\color{TitlePurple}\nouppercase{\leftmark}}
\fancyfoot[C]{\small\color{gray}— 第 \thepage\ 页 —}
\fancyfoot[R]{\small\color{gray}UCAS 22408}
\renewcommand{\headrulewidth}{0.5pt}

\pagenumbering{Roman}
\tableofcontents
\newpage
\pagenumbering{arabic}


\begin{targetbox}[本章复习目标与核心知识图谱]
\begin{itemize}
  \item 目的[1]:攻克非均匀密度介质与复杂异形旋转体容器抽水做功的高阶变力模型，掌握积分上下限参数化与变限微积分技巧
\end{itemize}
\end{targetbox}

\begin{targetbox}[本章复习目标与核心知识图谱]
\begin{itemize}
  \item 目的[2]:深入解析半无限长细棒、无限大平面薄板对空间质点引力矢量的反常积分模型与收敛性，领悟对称性相消与角度代换在空间引力场中的奇迹
\end{itemize}
\end{targetbox}

\begin{targetbox}[本章复习目标与核心知识图谱]
\begin{itemize}
  \item 目的[3]:攻克铅直平板静水压力中心（Center of Pressure）偏离形心的物理机制与截面二阶惯性矩（平行移轴定理）解析，掌握流体静力学压心严格位置公式
\end{itemize}
\end{targetbox}

\begin{targetbox}[本章复习目标与核心知识图谱]
\begin{itemize}
  \item 目的[4]:系统掌握空间刚体转动惯量微元与平行轴定理（Huygens-Steiner 定理）的微积分严格证明，打通一元积分与二重积分惯性矩的内在联系
\end{itemize}
\end{targetbox}

\begin{targetbox}[本章复习目标与核心知识图谱]
\begin{itemize}
  \item 目的[5]:攻克古鲁金第二定理（Pappus-Guldin 旋转体体积定理）与第一定理的联合运用，建立“形心 - 体积 - 侧面积 - 抽水功 - 静水压”五位一体秒杀网络
\end{itemize}
\end{targetbox}

\begin{targetbox}[本章复习目标与核心知识图谱]
\begin{itemize}
  \item 目的[6]:攻克带有需求价格弹性的变弹性利润最大化微分方程模型与垄断市场中税收无谓损失（Deadweight Loss）与社会总剩余最优化分析
\end{itemize}
\end{targetbox}


\section{1.非均匀密度介质与复杂异形容器的高阶抽水变力做功}

\begin{thmbox}[模型[1]: 密度随深度线性变化的盐水池抽水做功微元模型]

在实际工程与现代考研高阶大题中，液体介质往往不是理想均质常数密度，而是随水深呈现分层变化（如深水盐度分层、温度分层）。

\textbf{[1]} 非均匀密度介质抽水功的微元建立：
设水池深度为 $H$。池中液体的密度随深度 $x$ 呈线性分布：


\begin{equation*}
\rho(x) = \rho_0(1 + \beta x) \quad (\beta > 0)
\end{equation*}


其中 $\rho_0$ 为表层液体密度，$\beta$ 为密度梯度系数。重力加速度为 $g$。

\textbf{(1)}  \textbf{微元重力与提升位移}：

在深度 $x \in [0, H]$ 处取厚度为 $\mathrm{d}x$ 的水平薄层，截面积为 $A(x)$。

该薄层的微元质量为：


\begin{equation*}
\mathrm{d}m = \rho(x) A(x)\,\mathrm{d}x = \rho_0(1 + \beta x) A(x)\,\mathrm{d}x
\end{equation*}


薄层重力微元为：


\begin{equation*}
\mathrm{d}F = g\,\mathrm{d}m = \rho_0 g(1 + \beta x) A(x)\,\mathrm{d}x
\end{equation*}


将其提升至池顶面（$x = 0$）的垂直距离为 $h(x) = x$。

\textbf{(2)}  \textbf{做功微元与全池做功积分公式}：


\begin{equation*}
\mathrm{d}W = \mathrm{d}F \cdot x = \rho_0 g x(1 + \beta x) A(x)\,\mathrm{d}x
\end{equation*}


将整池液体完全抽出所做的总功为：


\begin{equation*}
W = \rho_0 g \int_0^H (x + \beta x^2) A(x)\,\mathrm{d}x
\end{equation*}


相比均质液体，非均匀介质额外引入了二阶矩项 $\int_0^H x^2 A(x)\,\mathrm{d}x$。

\textbf{[2]} 旋转抛物面异形水槽抽水功实例：
设一水槽由旋转抛物面 $z = c(x^2 + y^2)$ ($c > 0$) 绕 $z$ 轴对称形成，水槽总深为 $H$（槽底在 $z = 0$，槽口在 $z = H$）。槽内装满均质密度为 $\rho$ 的水。

以槽口平面为基准，向下建立水深坐标 $u = H - z$ ($u \in [0, H]$)。

在高度 $z$（即水深 $u = H - z$）处，水平截面为一个圆，其半径满足 $r^2(z) = \frac{z}{c} = \frac{H - u}{c}$。

水平截面积为：


\begin{equation*}
A(u) = \pi r^2 = \frac{\pi}{c}(H - u)
\end{equation*}


水深 $u$ 处的抽水做功微元为：


\begin{equation*}
\mathrm{d}W = \rho g u A(u)\,\mathrm{d}u = \frac{\pi \rho g}{c} u(H - u)\,\mathrm{d}u
\end{equation*}


全槽水抽空所做的总功为：


\begin{equation*}
W = \frac{\pi \rho g}{c} \int_0^H (H u - u^2)\,\mathrm{d}u = \frac{\pi \rho g}{c} \left[\frac{H u^2}{2} - \frac{u^3}{3}\right]_0^H = \frac{\pi \rho g H^3}{6c}
\end{equation*}


\end{thmbox}

\section{2.反常万有引力场与连续介质反常引力积分}

\begin{thmbox}[模型[3]: 半无限长均质细棒对空间任意质点的引力矢量反常积分与 $45^\circ$ 对称奇迹]

在反常积分物理应用中，半无限长连续介质的引力场呈现出极富对称美感的经典数学性质。

\textbf{[1]} 物理模型与引力矢量分解：
设有一根半无限长的均质细棒，放置在 $x$ 轴的正半轴 $[0, +\infty)$ 上，其线密度为常数 $\mu$。

在 $y$ 轴正半轴上点 $P(0, d)$ ($d > 0$) 处放置一质量为 $m$ 的质点。引力常数为 $G$。

在细棒上任意点 $x \in [0, +\infty)$ 处取微元段 $\mathrm{d}x$，质量为 $\mathrm{d}m_x = \mu\,\mathrm{d}x$。

该微元到质点 $P(0, d)$ 的距离为 $r = \sqrt{x^2 + d^2}$。

质元对质点 $P$ 产生的万有引力大小为：


\begin{equation*}
\mathrm{d}F = \frac{G m \mu}{r^2}\,\mathrm{d}x = \frac{G m \mu}{x^2 + d^2}\,\mathrm{d}x
\end{equation*}


引力方向由 $P(0, d)$ 指向 $(x, 0)$。引力矢量在 $x$ 轴与 $y$ 轴方向的分量分别为：


\begin{equation*}
\mathrm{d}F_x = \mathrm{d}F \cos\phi = \mathrm{d}F \frac{x}{\sqrt{x^2 + d^2}} = \frac{G m \mu x}{(x^2 + d^2)^{3/2}}\,\mathrm{d}x
\end{equation*}



\begin{equation*}
\mathrm{d}F_y = \mathrm{d}F \sin\phi = \mathrm{d}F \frac{d}{\sqrt{x^2 + d^2}} = \frac{G m \mu d}{(x^2 + d^2)^{3/2}}\,\mathrm{d}x
\end{equation*}


\textbf{[2]} 反常积分严格计算：
\textbf{(1)}  \textit{}水平分量 $F_x$（指向 $x$ 轴正向）\textit{}：


\begin{equation*}
F_x = G m \mu \int_0^{+\infty} \frac{x}{(x^2 + d^2)^{3/2}}\,\mathrm{d}x = \frac{1}{2} G m \mu \int_0^{+\infty} (x^2 + d^2)^{-3/2}\,\mathrm{d}(x^2 + d^2)
\end{equation*}



\begin{equation*}
= \frac{1}{2} G m \mu \left[-2(x^2 + d^2)^{-1/2}\right]_0^{+\infty} = G m \mu \left(0 - \left(-\frac{1}{d}\right)\right) = \frac{G m \mu}{d}
\end{equation*}


\textbf{(2)}  \textit{}垂直分量 $F_y$（指向 $y$ 轴负向）\textit{}：

作三角代换 $x = d \tan\theta$，$\mathrm{d}x = d \sec^2\theta\,\mathrm{d}\theta$。当 $x = 0$ 时 $\theta = 0$；当 $x \to +\infty$ 时 $\theta \to \frac{\pi}{2}$：


\begin{equation*}
F_y = G m \mu d \int_0^{\pi/2} \frac{d \sec^2\theta}{(d^2 \sec^2\theta)^{3/2}}\,\mathrm{d}\theta = \frac{G m \mu}{d} \int_0^{\pi/2} \cos\theta\,\mathrm{d}\theta = \frac{G m \mu}{d} [\sin\theta]_0^{\pi/2} = \frac{G m \mu}{d}
\end{equation*}


\textbf{[3]} 核心结论与几何奇迹：
比较两正交分量，惊人地发现：


\begin{equation*}
F_x = F_y = \frac{G m \mu}{d}
\end{equation*}


合引力大小为：


\begin{equation*}
F = \sqrt{F_x^2 + F_y^2} = \frac{\sqrt{2}G m \mu}{d}
\end{equation*}


合引力的方向角满足 $\tan\alpha = \frac{F_y}{F_x} = 1 \implies \alpha = 45^\circ$。

结论：\textit{}无论质点到细棒端点的距离 $d$ 为多大，半无限长细棒对质点的合万有引力方向恒定与坐标轴夹角为 $45^\circ$！\textit{} 这一反常积分性质在全国硕士研究生入学考试中具有极高的理论美感与考察价值。

\end{thmbox}
\begin{thmbox}[模型[4]: 无限大均质平面薄板对轴线上质点的反常引力与恒定场]

\textbf{[1]} 同心薄圆环微元法：
设有一块面密度为常数 $\sigma$ 的均质无限大平面薄板，占据整个 $xOy$ 平面。在垂直于平面的 $z$ 轴上点 $(0, 0, z)$ ($z > 0$) 处有一质量为 $m$ 的质点。

在平面上以原点为圆心，取半径为 $r$、宽度为 $\mathrm{d}r$ 的同心圆环。

圆环微元面积为 $\mathrm{d}S = 2\pi r\,\mathrm{d}r$，质量为 $\mathrm{d}m_r = 2\pi \sigma r\,\mathrm{d}r$。

圆环上各点到质点的距离均为 $R_0 = \sqrt{r^2 + z^2}$。

由圆周对称性，圆环对质点引力的水平分量在所有方向上严格抵消，合引力完全指向原点（沿 $z$ 轴负向）。

圆环对质点引力的垂直分量微元为：


\begin{equation*}
\mathrm{d}F_z = \frac{G m\,\mathrm{d}m_r}{r^2 + z^2} \cos\psi = \frac{G m (2\pi \sigma r\,\mathrm{d}r)}{r^2 + z^2} \frac{z}{\sqrt{r^2 + z^2}} = 2\pi G m \sigma z \frac{r\,\mathrm{d}r}{(r^2 + z^2)^{3/2}}
\end{equation*}


\textbf{[2]} 反常积分推导恒定场：
对于半径为 $R$ 的有限圆盘，轴向引力为：


\begin{equation*}
F_z(R) = 2\pi G m \sigma z \int_0^R \frac{r\,\mathrm{d}r}{(r^2 + z^2)^{3/2}} = 2\pi G m \sigma z \left[-\frac{1}{\sqrt{r^2 + z^2}}\right]_0^R = 2\pi G m \sigma \left(1 - \frac{z}{\sqrt{R^2 + z^2}}\right)
\end{equation*}


当圆盘半径趋向于无穷大（$R \to +\infty$）时：


\begin{equation*}
F_z = \lim_{R \to +\infty} 2\pi G m \sigma \left(1 - \frac{z}{\sqrt{R^2 + z^2}}\right) = 2\pi G m \sigma
\end{equation*}


结论：\textit{}无限大均质平板产生的引力场是一个严格的匀强引力场，其引力大小与质点到平板的垂直距离 $z$ 完全无关！\textit{} 这与电磁学中无限大均匀带电平板的匀强电场 $E = \frac{\sigma}{2\varepsilon_0}$ 具有完全同构的微积分数学本质。

\end{thmbox}

\section{3.液体压力中心（压力作用点）偏离形心的二阶矩模型}

\begin{thmbox}[模型[5]: 压力中心坐标公式与截面惯性矩（平行移轴定理渗透）]

在基础部分我们证明了单侧总压力等于形心处压强乘以面积 $P = p(\bar{y}) \cdot S$。但液体对平板各处的推力并不均匀（越深压强越大），因此\textbf{合水压力的等效作用点（压力中心 Center of Pressure）并不位于形心，而是必然位于形心下方}。

\textbf{[1]} 压力中心的力矩平衡推导：
设一铅直平板放置在水中，平板宽度函数为 $w(y)$，水深范围为 $y \in [a, b]$。

在深度 $y$ 处的微元压力为 $\mathrm{d}P = \rho g y w(y)\,\mathrm{d}y$。

该微元压力对水面交线（$x$ 轴）的转动力矩微元为：


\begin{equation*}
\mathrm{d}M_x = y \cdot \mathrm{d}P = \rho g y^2 w(y)\,\mathrm{d}y
\end{equation*}


总水压力 $P$ 对水面轴的总力矩为：


\begin{equation*}
M_x = \int_a^b \mathrm{d}M_x = \rho g \int_a^b y^2 w(y)\,\mathrm{d}y
\end{equation*}


设合力 $P$ 的集中等效作用点深度为 $y_p$（压力中心深度）。根据力矩等效原理：


\begin{equation*}
P \cdot y_p = M_x \implies y_p = \frac{M_x}{P} = \frac{\rho g \int_a^b y^2 w(y)\,\mathrm{d}y}{\rho g \int_a^b y w(y)\,\mathrm{d}y} = \frac{\int_a^b y^2 w(y)\,\mathrm{d}y}{\int_a^b y w(y)\,\mathrm{d}y}
\end{equation*}


\textbf{[2]} 与截面惯性矩与平行移轴定理的深刻联系：
在截面几何学中：

\textbf{(1)}  \textbf{一阶静矩}：$S_x = \int_a^b y w(y)\,\mathrm{d}y = \bar{y} \cdot S$（$\bar{y}$ 为形心深度，$S$ 为平板总面积）；

\textbf{(2)}  \textbf{二阶惯性矩}：$I_x = \int_a^b y^2 w(y)\,\mathrm{d}y$ 为平板对水面轴的转动惯性矩。

根据力学与微积分的\textbf{平行移轴定理}：


\begin{equation*}
I_x = I_{\bar{x}} + \bar{y}^2 S
\end{equation*}


其中 $I_{\bar{x}}$ 为平板绕通过自身几何形心且平行于水面的轴的自惯性矩（截面二次轴矩）。

将此关系代入压力中心深度表达式：


\begin{equation*}
y_p = \frac{I_x}{S_x} = \frac{I_{\bar{x}} + \bar{y}^2 S}{\bar{y} S} = \bar{y} + \frac{I_{\bar{x}}}{\bar{y} S}
\end{equation*}


\textbf{[3]} 压力中心三大核心物理推论：
\textbf{(1)}  \textbf{压心必在形心下方}：由于任何平面图形的二次轴矩 $I_{\bar{x}} > 0$ 且面积 $S > 0$，故 $\frac{I_{\bar{x}}}{\bar{y} S} > 0$，因此 $y_p > \bar{y}$ 恒成立！

\textbf{(2)}  \textbf{渐近逼近性}：当平板浸没非常深（即 $\bar{y} \to +\infty$）时，$\frac{I_{\bar{x}}}{\bar{y} S} \to 0$，压力中心逐渐与形心重合。

\textbf{(3)}  \textbf{矩形水闸平板特例}：

若矩形平板高为 $h$、宽为 $b$，顶边在水面（$a = 0, b = h$）。

其形心为 $\bar{y} = \frac{h}{2}$，自身惯性矩为 $I_{\bar{x}} = \frac{b h^3}{12}$，面积 $S = b h$。

代入公式：


\begin{equation*}
y_p = \frac{h}{2} + \frac{\frac{b h^3}{12}}{\frac{h}{2} \cdot b h} = \frac{h}{2} + \frac{h}{6} = \frac{2}{3}h
\end{equation*}


结论：\textit{}顶边与水面齐平的垂直矩形闸门，其水压力等效作用点恰好位于水深 $\frac{2}{3}h$ 处！\textit{}

\end{thmbox}

\section{4.空间刚体转动惯量与平行轴定理的微积分严格证明}

\begin{thmbox}[定理[1]: 平行轴定理（Huygens-Steiner 定理）的一般微积分证明]

\textbf{[1]} 定理叙述：
刚体对任一转轴 $z$ 的转动惯量 $I_z$，等于刚体对通过其质心且与该轴平行的转轴 $z_C$ 的转动惯量 $I_C$，加上刚体总质量 $M$ 与两轴垂直距离 $d$ 的平方之乘积：


\begin{equation*}
I_z = I_C + M d^2
\end{equation*}


\textbf{[2]} 严格定积分证明：
建立笛卡尔直角坐标系，设质心轴 $z_C$ 穿过坐标原点 $(0, 0)$，平行轴 $z$ 穿过点 $(x_0, y_0)$。

两平行轴之间的垂直几何距离为 $d = \sqrt{x_0^2 + y_0^2} \implies d^2 = x_0^2 + y_0^2$。

设刚体内任意微元 $\mathrm{d}m$ 位于点 $(x, y, z)$。

微元到质心轴 $z_C$ 的垂直距离平方为 $r_C^2 = x^2 + y^2$。

微元到平行轴 $z$ 的垂直距离平方为：


\begin{equation*}
r^2 = (x - x_0)^2 + (y - y_0)^2 = (x^2 + y^2) + (x_0^2 + y_0^2) - 2(x_0 x + y_0 y) = r_C^2 + d^2 - 2(x_0 x + y_0 y)
\end{equation*}


对刚体全域进行积分：


\begin{equation*}
I_z = \int r^2\,\mathrm{d}m = \int r_C^2\,\mathrm{d}m + d^2 \int \mathrm{d}m - 2x_0 \int x\,\mathrm{d}m - 2y_0 \int y\,\mathrm{d}m
\end{equation*}


注意到：

\textbf{(1)}  $\int r_C^2\,\mathrm{d}m = I_C$（绕质心轴的转动惯量）；

\textbf{(2)}  $\int \mathrm{d}m = M$（刚体总质量）；

\textbf{(3)}  由于原点选在质心位置，根据质心坐标定义：


\begin{equation*}
\bar{x} = \frac{\int x\,\mathrm{d}m}{M} = 0 \implies \int x\,\mathrm{d}m = 0
\end{equation*}



\begin{equation*}
\bar{y} = \frac{\int y\,\mathrm{d}m}{M} = 0 \implies \int y\,\mathrm{d}m = 0
\end{equation*}


后两项交叉项严格恒等于零！

因此：


\begin{equation*}
I_z = I_C + M d^2
\end{equation*}


证毕。

\end{thmbox}

\section{5.古鲁金第二定理与三位一体秒杀体系}

\begin{thmbox}[定理[2]: 古鲁金第二定理（Pappus-Guldin 旋转体体积定理）推导与联合运用]

在基础篇我们系统掌握了关于侧面积的古鲁金第一定理（$A = 2\pi \bar{r} l$）。而古鲁金第二定理则是关于平面图形旋转体体积的孪生定理。

\textbf{[1]} 定理表述：
平面上一有界闭区域 $D$ 绕同平面内一条与该区域不相交（至多在边界相交）的直线轴 $L_0$ 旋转一周所形成的旋转体的体积 $V$，等于该平面区域的面积 $S$ 与该区域形心 $(\bar{x}, \bar{y})$ 绕旋转轴旋转一周所经过的圆周路程的乘积：


\begin{equation*}
V = 2\pi r(\bar{x}, \bar{y}) \cdot S
\end{equation*}


其中 $r(\bar{x}, \bar{y})$ 为平面区域形心到旋转轴 $L_0$ 的垂直距离。

\textbf{[2]} 严格微积分推导（以绕纵轴为例）：
设区域 $D$ 位于右半平面，由连续曲线 $y_1(x) \leqslant y \leqslant y_2(x)$ ($a \leqslant x \leqslant b$, $a \geqslant 0$) 围成。

采用柱壳法（Cylindrical Shells Method），旋转体体积微元为：


\begin{equation*}
\mathrm{d}V = 2\pi x [y_2(x) - y_1(x)]\,\mathrm{d}x
\end{equation*}


总旋转体体积为：


\begin{equation*}
V = 2\pi \int_a^b x [y_2(x) - y_1(x)]\,\mathrm{d}x = 2\pi M_y
\end{equation*}


其中 $M_y = \iint_D x\,\mathrm{d}x\mathrm{d}y = \int_a^b x [y_2(x) - y_1(x)]\,\mathrm{d}x$ 为平面区域 $D$ 对 $y$ 轴的一阶静矩。

由平面区域形心坐标公式 $\bar{x} = \frac{M_y}{S} \implies M_y = \bar{x} \cdot S$。

代入即得：


\begin{equation*}
V = 2\pi \bar{x} \cdot S
\end{equation*}


证毕。

\textbf{[3]} 圆环体（Torus）的古鲁金双定理秒杀大典：
设圆盘 $(x - b)^2 + y^2 \leqslant R^2$ ($b > R$) 绕 $y$ 轴旋转一周形成圆环体。

\textbf{(1)}  \textbf{表面积（古鲁金第一定理）}：

母线圆周长 $l = 2\pi R$，形心为圆心 $(b, 0)$，到轴距离 $r = b$：


\begin{equation*}
A = 2\pi r \cdot l = 2\pi b \cdot 2\pi R = 4\pi^2 b R
\end{equation*}


\textbf{(2)}  \textbf{体积（古鲁金第二定理）}：

平面圆面积 $S = \pi R^2$，形心到轴距离 $r = b$：


\begin{equation*}
V = 2\pi r \cdot S = 2\pi b \cdot \pi R^2 = 2\pi^2 b R^2
\end{equation*}


无需任何复杂的参数代换或二重积分，两行公式直接写出答案。

\end{thmbox}

\section{6.经济学深水区：需求弹性微分方程与社会福利优化}

\begin{thmbox}[模型[6]: 需求价格弹性 $\eta(P)$ 与边际收益方程的联合积分]

\textbf{[1]} 需求价格弹性定义与符号约定：
需求价格弹性 $\eta$ 定义为需求量变动的相对百分比与价格变动相对百分比之比的绝对值：


\begin{equation*}
\eta = -\frac{\mathrm{d}Q / Q}{\mathrm{d}P / P} = -\frac{P}{Q}\frac{\mathrm{d}Q}{\mathrm{d}P}
\end{equation*}


\textbf{[2]} 阿莫罗索-罗宾逊（Amoroso-Robinson）公式的微分推导：
总收益为价格与销售量的乘积：$R = P \cdot Q(P)$。

总收益对价格 $P$ 求导：


\begin{equation*}
\frac{\mathrm{d}R}{\mathrm{d}P} = Q + P \frac{\mathrm{d}Q}{\mathrm{d}P} = Q \left(1 + \frac{P}{Q}\frac{\mathrm{d}Q}{\mathrm{d}P}\right) = Q (1 - \eta)
\end{equation*}


边际收益 $MR = \frac{\mathrm{d}R}{\mathrm{d}Q}$ 与价格 $P$ 的关系为：


\begin{equation*}
MR = \frac{\mathrm{d}R}{\mathrm{d}Q} = \frac{\mathrm{d}R / \mathrm{d}P}{\mathrm{d}Q / \mathrm{d}P} = \frac{Q(1 - \eta)}{\mathrm{d}Q / \mathrm{d}P} = \frac{Q}{\mathrm{d}Q / \mathrm{d}P} (1 - \eta) = -\frac{P}{\eta}(1 - \eta) = P \left(1 - \frac{1}{\eta}\right)
\end{equation*}


经济学深刻结论：

\textbf{(1)}  若 $\eta > 1$（富有弹性），则 $MR > 0$，降价可增加总收益；

\textbf{(2)}  若 $\eta < 1$（缺乏弹性），则 $MR < 0$，降价会导致总收益减少；

\textbf{(3)}  若 $\eta = 1$（单位弹性），则 $MR = 0$，总收益达到极大值。

\textbf{[3]} 由恒弹性反求需求函数与福利剩余微分方程：
若已知某产品在整个消费区间上具有常数需求弹性 $\eta = k$ ($k > 0, k \neq 1$)，且当价格为 $P_0$ 时销售量为 $Q_0$。

可列出可分离变量微分方程：


\begin{equation*}
-\frac{P}{Q}\frac{\mathrm{d}Q}{\mathrm{d}P} = k \implies \frac{\mathrm{d}Q}{Q} = -k \frac{\mathrm{d}P}{P}
\end{equation*}


积分得：


\begin{equation*}
\ln Q = -k \ln P + C_1 \implies Q(P) = A P^{-k}
\end{equation*}


代入初始条件 $Q(P_0) = Q_0$，得需求函数解析式：


\begin{equation*}
Q = Q_0 \left(\frac{P}{P_0}\right)^{-k} \iff P(Q) = P_0 \left(\frac{Q}{Q_0}\right)^{-\frac{1}{k}}
\end{equation*}


当 $k > 1$ 时，反常积分求消费者剩余收敛：


\begin{equation*}
CS = \int_0^{Q_0} [P(Q) - P_0]\,\mathrm{d}Q = P_0 \int_0^{Q_0} \left[\left(\frac{Q}{Q_0}\right)^{-\frac{1}{k}} - 1\right]\,\mathrm{d}Q = \frac{P_0 Q_0}{k - 1}
\end{equation*}


\end{thmbox}
\begin{thmbox}[模型[7]: 税收无谓损失（Deadweight Loss）与剩余面积极值优化]

\textbf{[1]} 征税前后的市场均衡位移：
设某商品的线性需求函数为 $P = a - b Q$ ($a, b > 0$)，供给函数为 $P = c + d Q$ ($c > 0, d > 0, a > c$)。

未征税前均衡产量与价格为：


\begin{equation*}
Q_0 = \frac{a - c}{b + d}, \quad P_0 = \frac{a d + b c}{b + d}
\end{equation*}


政府对每单位商品征收单位税额为 $t$ 的从量税。此时供给曲线向上平移 $t$：


\begin{equation*}
P_S^t(Q) = c + d Q + t
\end{equation*}


新的供需均衡由 $a - b Q_t = c + d Q_t + t$ 确定：


\begin{equation*}
Q_t = \frac{a - c - t}{b + d} = Q_0 - \frac{t}{b + d}
\end{equation*}


产量的减少量为：


\begin{equation*}
\Delta Q = Q_0 - Q_t = \frac{t}{b + d}
\end{equation*}


\textbf{[2]} 社会总福利的无谓损失（哈伯格三角形 Deadweight Triangle）：
征税后：

\textbf{(1)}  消费者剩余减少量 $\Delta CS$ 对应梯形面积；

\textbf{(2)}  生产者剩余减少量 $\Delta PS$ 对应梯形面积；

\textbf{(3)}  政府获得税收总额为 $T = t \cdot Q_t$。

政府税收收入不足以弥补消费者与生产者损失之和，由此产生社会总福利的净损失（无谓损失 $DWL$）。

根据定积分面积差计算，该无谓损失恰好为两线之间夹成的三角形微积分面积：


\begin{equation*}
DWL = \int_{Q_t}^{Q_0} [P_D(Q) - P_S(Q)]\,\mathrm{d}Q = \int_{Q_t}^{Q_0} [(a - b Q) - (c + d Q)]\,\mathrm{d}Q
\end{equation*}



\begin{equation*}
= \frac{1}{2} [(P_D(Q_t) - P_S(Q_t)) + (P_D(Q_0) - P_S(Q_0))] \cdot (Q_0 - Q_t)
\end{equation*}


由于在 $Q_t$ 处买卖双方价差恰为税额 $t$，而在 $Q_0$ 处价差为 $0$：


\begin{equation*}
DWL = \frac{1}{2} (t + 0) \cdot \Delta Q = \frac{1}{2} t \cdot \frac{t}{b + d} = \frac{t^2}{2(b + d)}
\end{equation*}


重要结论：\textit{}税收无谓损失与单位税额的平方 $t^2$ 成正比！\textit{} 这一结论是公共财政学与福利经济学中最著名的微积分极值法则。
\end{thmbox}

\end{document}
